\section{Crystals in quasi-coherent modules}
\label{section-quasi-coherent-crystals}

\noindent
In Situation \ref{situation-affine}.
Set $X = \Spec(C)$ and $S = \Spec(A)$. We are going to
classify crystals in quasi-coherent modules on $\text{Cris}(X/S)$.
Before we do so we fix some notation.

\medskip\noindent
Choose a polynomial ring $P = A[x_i]$ over $A$ and a surjection $P \to C$
of $A$-algebras with kernel $J = \Ker(P \to C)$. Set
\begin{equation}
\label{equation-D}
D = \lim_e D_{P, \gamma}(J) / p^eD_{P, \gamma}(J)
\end{equation}
for the $p$-adically completed divided power envelope.
This ring comes with a divided power ideal $\bar J$ and divided power
structure $\bar \gamma$, see Lemma \ref{lemma-list-properties}.
Set $D_e = D/p^eD$ and denote $\bar J_e$ the image of $\bar J$ in $D_e$.
We will use the short hand
\begin{equation}
\label{equation-omega-D}
\Omega_D = \lim_e \Omega_{D_e/A, \bar\gamma} =
\lim_e \Omega_{D/A, \bar\gamma}/p^e\Omega_{D/A, \bar\gamma}
\end{equation}
for the $p$-adic completion of the module of divided power differentials,
see Lemma \ref{lemma-differentials-completion}.
It is also the $p$-adic completion of
$\Omega_{D_{P, \gamma}(J)/A, \bar\gamma}$
which is free on $\text{d}x_i$, see
Lemma \ref{lemma-module-differentials-divided-power-envelope}.
Hence any element of $\Omega_D$ can be written uniquely as a sum
$\sum f_i\text{d}x_i$ with for all $e$ only finitely many $f_i$
not in $p^eD$. Moreover, the maps
$\text{d}_{D_e/A, \bar\gamma} : D_e \to \Omega_{D_e/A, \bar\gamma}$
fit together to define a divided power $A$-derivation
\begin{equation}
\label{equation-derivation-D}
\text{d} : D \longrightarrow \Omega_D
\end{equation}
on $p$-adic completions.

\medskip\noindent
We will also need the ``products $\Spec(D(n))$ of $\Spec(D)$'', see
Proposition \ref{proposition-compute-cohomology} and its proof for an
explanation. Formally these are defined as follows. For $n \geq 0$ let
$J(n) = \Ker(P \otimes_A \ldots \otimes_A P \to C)$ where
the tensor product has $n + 1$ factors. We set
\begin{equation}
\label{equation-Dn}
D(n) = \lim_e
D_{P \otimes_A \ldots \otimes_A P, \gamma}(J(n))/
p^eD_{P \otimes_A \ldots \otimes_A P, \gamma}(J(n))
\end{equation}
equal to the $p$-adic completion of the divided power envelope.
We denote $\bar J(n)$ its divided power ideal and $\bar \gamma(n)$
its divided powers. We also introduce $D(n)_e = D(n)/p^eD(n)$ as well
as the $p$-adically completed module of differentials
\begin{equation}
\label{equation-omega-Dn}
\Omega_{D(n)} = \lim_e \Omega_{D(n)_e/A, \bar\gamma} =
\lim_e \Omega_{D(n)/A, \bar\gamma}/p^e\Omega_{D(n)/A, \bar\gamma}
\end{equation}
and derivation
\begin{equation}
\label{equation-derivation-Dn}
\text{d} : D(n) \longrightarrow \Omega_{D(n)}
\end{equation}
Of course we have $D = D(0)$. Note that the rings $D(0), D(1), D(2), \ldots$
form a cosimplicial object in the category of divided power rings.

